\bookchapter{Capitolo 0: Ehi, prima di iniziare}

Prima di esplorare ufficialmente, parliamo di alcune regole di base per assicurarci che tu capisca tutto ciò che diremo dopo. Non aver fretta, affilare la scure non ritarda il taglio della legna.

\booksection{Cosa può fare questo libro per te? (E cosa non può)}

Può aiutarti a capire sistematicamente come usare il cellulare, a capire la logica dietro le funzioni; inoltre, che tu sia del team Android o Apple, la maggior parte delle idee è universale. Ti guiderà in esercizi pratici, facendoti passare da “saper usare” a “capire come usare”.

Non è il manuale di un modello specifico di cellulare, né notizie tecnologiche sull'inseguire le ultime novità, e tanto meno una guida per farti assorbire dal cellulare. A dire il vero, se ti aspetti “come scrollare video in modo più piacevole”, potresti restare deluso — ci preoccupiamo più di come far servire davvero il cellulare per te, non il contrario.

\booksection{Spiegazione dei nostri “codici”}

Per spiegare meglio, il libro usa alcuni piccoli segni; questi segni sono come un'intesa tacita tra noi:

\begin{itemize}
  \item \textbf{[Nome pulsante]}: rappresenta un pulsante specifico sullo schermo o il nome di un'app, ad esempio [Impostazioni], [OK].
  \item \textbf{>}: rappresenta l'ordine delle operazioni, ad esempio "tocca [Impostazioni] > [Suono]" significa prima toccare Impostazioni, poi Suono.
  \item \textbf{Grassetto}: molto importante! Potrebbe essere un concetto chiave o un avvertimento di sicurezza, prestaci attenzione.
  \item \emph{Corsivo}: piccoli pensieri integrativi o metafore per aiutarti a capire meglio.
\end{itemize}

\booksection{Android e iOS, parliamo di entrambi}

Indovina un po'? Alcune cose hanno nomi diversi ma fanno la stessa cosa. Ad esempio: applicazione (Android) e App (iOS), Impostazioni (entrambi le chiamano così), Gestione file (Android) e File (iOS). Basta sapere che si tratta della stessa cosa, non farti confondere dai termini.

Inoltre, ora nel campo Android c'è anche il sistema HarmonyOS di Huawei, che ha le sue caratteristiche; ne parleremo in modo specifico più avanti. In sintesi, qualsiasi cellulare tu usi, questo libro può aiutarti.

\booksection{Capacità del cellulare: perché ho comprato 128G ma lo spazio disponibile è meno di 110G?}

Le unità di capacità sono note: 1TB = 1024GB, 1GB = 1024MB, 1MB = 1024KB... ma c'è una trappola: i produttori di dischi rigidi di solito usano base 1000, mentre i sistemi usano base 1024, quindi il disco “128G” che hai comprato “si restringe” nella capacità effettiva.

Ma i motivi principali del “restringimento” della capacità del cellulare sono tre:

\begin{enumerate}
  \item \textbf{Il sistema stesso occupa spazio}: il sistema operativo è come i muri portanti di una casa, deve esserci. Oggi i sistemi occupano facilmente 10-20 GB, è normale.
  \item \textbf{App preinstallate dal produttore}: le app che arrivano con il nuovo telefono occupano spazio; alcune si possono disinstallare, altre no.
  \item \textbf{Perdite di partizionamento e formattazione}: motivi tecnici, come la fodera dello zaino che occupa volume.
\end{enumerate}

\begin{tcbinfo}{Piccolo consiglio per l'acquisto}
Oggi 128 GB sono ormai insufficienti; partire da 256 GB è più prudente. Se ti piace scattare foto o giocare, valuta direttamente 512 GB o più. Ricorda, per lo spazio di archiviazione, meglio abbondare che essere a corto.

\end{tcbinfo}

\begin{tcbpractice}{Fai pratica: guarda la "carta d'identità" del tuo cellulare}
Apri [Impostazioni], trova Info sul telefono (o opzione simile) e annota: modello, versione di sistema, spazio di archiviazione totale e residuo. Questo è il primo passo per conoscerlo. Già che ci sei, controlla lo stato di salute della batteria per avere un'idea.

\end{tcbpractice}

\booksection{Il tuo cellulare è una “libreria di asset digitali”}

Prima di iniziare con le conoscenze specifiche, stabiliamo una cognizione di base:\textbf{il cellulare non è più solo uno strumento di comunicazione, è la nostra libreria di asset digitali più importante}. Contatti, cronologia chat, foto di famiglia, documenti di lavoro, app bancarie, codici di pagamento, dati sanitari, documenti elettronici... questi dati messi insieme valgono spesso più del cellulare stesso.

\begin{tcbinfo}{Perché è importante?}
Se il cellulare viene perso o si guasta, perdi non solo qualche migliaio di euro di hardware,\textbf{ma probabilmente anni di dati non riproducibili}— le foto della crescita dei figli, le voci di parenti scomparsi, contratti importanti senza backup. Abitudini di organizzazione e protezione scientifiche possono proteggerti su tre dimensioni:

\begin{itemize}
  \item \textbf{Maggiore efficienza}: una classificazione dei file regolamentata rende più veloce la ricerca, il cellulare non rallenta nel tempo.
  \item \textbf{Vita più lunga}: una gestione ragionevole dell'archiviazione riduce il carico di lettura/scrittura dei chip, la batteria e l'archiviazione durano di più.
  \item \textbf{Sicurezza dei dati}: gestione a livelli + backup periodico + restrizione dei permessi, allontana i "disastri dei dati".
\end{itemize}

\end{tcbinfo}

\booksubsection{La nostra filosofia di base: “Prevenzione prima, combinare prevenzione e cura”}

Come per la salute non si può aspettare di ammalarsi per andare dal medico,\textbf{anche la gestione del cellulare va affrontata a partire dalle “abitudini quotidiane”}. Questo libro non sottolinea una tecnica avanzata specifica, ma questi principi semplici:

\begin{itemize}
  \item \textbf{Prevenzione prima}: installare app solo dallo store ufficiale, rimuovere dati sensibili dalle foto prima di inviarle, non usare password duplicate — queste "piccole azioni" costano pochissimo ma bloccano la maggior parte dei rischi.
  \item \textbf{Controlli periodici}: riorganizzare la cartella download una volta al mese, audit dei permessi una volta per trimestre, backup completo una volta ogni semestre — un ritmo fisso è cento volte più efficace del soccorso d'emergenza.
  \item \textbf{Lasciarsi una via d'uscita}: attivare la sincronizzazione cloud, annotare account e password, scrivere un "inventario del patrimonio digitale" — per garantire che tu e la tua famiglia possiate gestire con calma in caso di imprevisti.
\end{itemize}
